\pagestyle{empty}
\tight
\begin{tblr}{width=\textwidth,colspec={|Q[c,m,1.9cm]|X[l,m]|},hlines,vlines,hspan=minimal,row{1}={bg=headblue},rowsep=1pt,colsep=3pt}
\SetCell{c,m}\hfil\logo\hfil & \textbf{POLITEKNIK NEGERI MALANG}\par\textbf{JURUSAN TEKNOLOGI INFORMASI}\par\textbf{PROGRAM STUDI: D4 TEKNIK INFORMATIKA} \\
\SetCell[c=2]{c} \textbf{RENCANA PEMBELAJARAN SEMESTER (RPS)} & \\
\end{tblr}
\begin{longtblr}{width=\textwidth,colspec={|Q[l,m,1.9cm]|X[0.85,l,m]|X[1.45,l,m]|X[0.75,l,m]|X[1.15,l,m]|},hlines,vlines,hspan=minimal,rowsep=1pt,colsep=2pt,row{1,3}={bg=headgray,font=\bfseries}}
MATA KULIAH & KODE & BOBOT (SKS)/jam & SEMESTER & TGL. PENYUSUNAN \\
Magang & RTI257001 & 20 SKS / 40 jam/minggu & 7 & 1 Juli 2025 \\
OTORISASI & Dosen Pengembang RPS & Koordinator RMK & \SetCell[c=2]{l} Ka PRODI &  \\
 & Triana Fatmawati, S.T., M.T. &  & \SetCell[c=2]{l}{\vspace{8mm}\par Dr. Ely Setyo Astuti, S.T., M.T.} &  \\
\SetCell[r=14]{l} Capaian Pembelajaran (CP) & \SetCell[c=4]{l,bg=headgray,font=\bfseries} Capaian Pembelajaran Lulusan Program Studi (CPL-Prodi) &  &  &  \\
 & CPL03 & \SetCell[c=3]{l} Mampu mengelola tim dan diri sendiri secara profesional serta mengomunikasikan dan mempresentasikan gagasan maupun hasil kerja secara efektif baik lisan maupun tulisan. &  &  \\
 & CPL06 & \SetCell[c=3]{l} Mampu merancang, mengimplementasikan, menguji, melakukan deployment, memelihara, serta menjamin mutu perangkat lunak yang menjawab permasalahan dan memenuhi kebutuhan pengguna. &  &  \\
 & CPL08 & \SetCell[c=3]{l} Mampu mengelola proyek teknologi informasi secara profesional melalui perencanaan, koordinasi, komunikasi, dan optimalisasi sumber daya. &  &  \\
 & \SetCell[c=4]{l,bg=headgray,font=\bfseries} Capaian Pembelajaran mata kuliah (CPL-MK) &  &  &  \\
 & CPMK0301 & \SetCell[c=3]{l} Mahasiswa mampu berkomunikasi secara profesional lisan dan tulisan dalam lingkungan kerja industri TI. &  &  \\
 & CPMK0303 & \SetCell[c=3]{l} Mahasiswa mampu mempresentasikan hasil kerja dan laporan magang secara efektif kepada pembimbing dan penguji. &  &  \\
 & CPMK0804 & \SetCell[c=3]{l} Mahasiswa mampu menyelesaikan proyek atau tugas industri secara mandiri dan profesional selama magang. &  &  \\
 & CPMK0607 & \SetCell[c=3]{l} Mahasiswa mampu mengaplikasikan kompetensi rekayasa perangkat lunak dalam konteks industri nyata. &  &  \\
 & \SetCell[c=4]{l,bg=headgray,font=\bfseries} Kemampuan akhir tiap tahapan belajar (Sub-CPMK) &  &  &  \\
 & SCPMK0301-05701 & \SetCell[c=3]{l} Mahasiswa mampu berkomunikasi profesional secara lisan dan tulisan selama kegiatan magang di industri TI. &  &  \\
 & SCPMK0303-05702 & \SetCell[c=3]{l} Mahasiswa mampu menyusun dan mempresentasikan laporan magang yang komprehensif kepada pembimbing akademik dan penguji. &  &  \\
 & SCPMK0804-05703 & \SetCell[c=3]{l} Mahasiswa mampu menyelesaikan target pekerjaan industri yang diberikan pembimbing lapangan secara mandiri dan tepat waktu. &  &  \\
 & SCPMK0607-05704 & \SetCell[c=3]{l} Mahasiswa mampu mengaplikasikan minimal satu kompetensi rekayasa perangkat lunak dalam proyek atau tugas industri nyata. &  &  \\
\SetCell[r=2]{l} Penilaian & \SetCell[c=4]{l} *Nilai CPMK didapat dari agregasi nilai SUB-CPMK yang dibebankan &  &  &  \\
 & \SetCell[c=4]{l}{\tiny\begin{tblr}{width=\linewidth,colspec={|Q[c,m,0.1\linewidth]|Q[c,m,0.16\linewidth]|X[l,m]|X[l,m]|X[l,m]|X[l,m]|Q[c,m,0.08\linewidth]|},hlines,vlines,hspan=minimal,rowsep=1pt,colsep=0.7pt,row{1,2}={bg=headgray,font=\bfseries}}
Kode CPMK & Kode SCPMK & \SetCell[c=4]{c} Bobot per Bentuk Penilaian & & & & Total Bobot \\
 & & Logbook & Nilai Lapangan & Laporan & Sidang & \\
CPMK0301 & SCPMK0301-05701 & 10\%: logbook komunikasi & & 8\%: laporan magang & & 18\% \\
CPMK0303 & SCPMK0303-05702 & & & 12\%: laporan magang & 15\%: sidang magang & 27\% \\
CPMK0804 & SCPMK0804-05703 & 10\%: logbook kegiatan & 15\%: nilai pembimbing lapangan & & & 25\% \\
CPMK0607 & SCPMK0607-05704 & & 15\%: nilai kompetensi RPL & & 15\%: sidang kompetensi & 30\% \\
\SetCell[c=6]{r}\textbf{Total Per Penilaian} & & & & & & \textbf{100\%} \\
\end{tblr}} &  &  &  \\
Diskripsi Singkat Mata Kuliah & \SetCell[c=4]{l} Magang adalah program pembelajaran luar kampus (MBKM) 20 SKS yang menempatkan mahasiswa di industri TI selama ±4 bulan untuk mengaplikasikan kompetensi akademis, membangun profesionalisme, dan menghasilkan laporan ilmiah berbasis pengalaman kerja nyata. &  &  &  \\
Bahan Kajian / Materi Pembelajaran & \SetCell[c=4]{l}\begin{enumerate}\item Pembekalan magang, etika profesi, dan orientasi industri\item Praktik kerja industri dan logbook kegiatan\item Aplikasi kompetensi rekayasa perangkat lunak di industri\item Penulisan laporan magang dan komunikasi profesional\item Presentasi dan sidang magang\end{enumerate} &  &  &  \\
\SetCell[r=2]{l} Pustaka & \SetCell[c=4]{l} \textbf{Utama:}\begin{enumerate}\item Buku Pedoman Magang D4 TI Polinema 2025.\item Panduan MBKM Kemdikbud 2020.\item Pedoman Penulisan Laporan Magang Polinema.\end{enumerate} &  &  &  \\
 & \SetCell[c=4]{l} \textbf{Pendukung:} Materi LMS, dokumen industri/pedoman magang, dan bahan bimbingan. &  &  &  \\
Media Pembelajaran & \SetCell[c=2]{l} \textbf{Software:} Aplikasi logbook digital, LMS, sistem pelaporan. &  & \SetCell[c=2]{l} \textbf{Hardware:} Komputer/laptop, perangkat kerja industri. &  \\
Nama Dosen Pengampu & \SetCell[c=4]{l} Tim Pengajar Program Studi D4 TI &  &  &  \\
Matakuliah Syarat & \SetCell[c=4]{l} Lulus Semester 1--6 (minimal 120 SKS) &  &  &  \\
\end{longtblr}
\begin{longtblr}{width=\textwidth,colspec={|Q[c,m,0.06\textwidth]|X[1.35,l,m]|X[1.08,l,m]|X[1.16,l,m]|X[0.7,l,m]|X[1.24,l,m]|X[1.06,l,m]|X[0.98,l,m]|Q[c,m,0.05\textwidth]|},hlines,vlines,hspan=minimal,rowhead=2,row{1,2}={bg=headgreen,font=\bfseries},rowsep=1pt,colsep=1.5pt}
Mgg.\par Ke & Kemampuan Akhir Yang Direncanakan (Sub-CP-MK) & Bahan kajian (Materi Pembelajaran) & Bentuk dan Metode Pembelajaran & Estimasi Waktu & Pengalaman Belajar Mahasiswa & Kriteria \& Bentuk Penilaian & Indikator Penilaian & Bobot Penilaian (\%) \\
(1) & (2) & (3) & (4) & (5) & (6) & (7) & (8) & (9) \\
1 & SCPMK0301-05701, SCPMK0804-05703 & Pembekalan magang: orientasi, etika kerja, dan kontrak belajar. & Modalitas: Praktik Lapangan dan Bimbingan. Metode: praktik industri, monitoring berkala, dan bimbingan pembimbing akademik/lapangan. & 1 x 2 x 50' tatap muka; persiapan mandiri. & Mahasiswa mengikuti sesi pembekalan, memahami prosedur magang, etika kerja industri, dan menandatangani kontrak belajar. & Bentuk penilaian: Kehadiran dan partisipasi pembekalan. Mengacu rubrik RTM. & Ketepatan waktu; kualitas dokumen; profesionalisme. & 2 \\
2 & SCPMK0804-05703 & Penempatan dan pengenalan lingkungan industri. & Modalitas: Praktik Lapangan dan Bimbingan. Metode: praktik industri, monitoring berkala, dan bimbingan pembimbing akademik/lapangan. & Disesuaikan jadwal magang. & Mahasiswa ditempatkan di perusahaan industri TI, mengenal lingkungan kerja, dan mulai memahami budaya organisasi. & Bentuk penilaian: Laporan penempatan awal. Mengacu rubrik RTM. & Ketepatan waktu; kualitas dokumen; profesionalisme. & 2 \\
3 & SCPMK0607-05704 & Penyerahan ke industri dan mulai kerja. & Modalitas: Praktik Lapangan dan Bimbingan. Metode: praktik industri, monitoring berkala, dan bimbingan pembimbing akademik/lapangan. & Disesuaikan jadwal magang. & Mahasiswa diserahkan ke pembimbing lapangan, mendapatkan tugas awal, dan mulai berkontribusi dalam proyek/pekerjaan industri. & Bentuk penilaian: Logbook awal minggu pertama. Mengacu rubrik RTM. & Ketepatan waktu; kualitas dokumen; profesionalisme. & 2 \\
4 & SCPMK0301-05701, SCPMK0804-05703 & Monitoring bulan 1: logbook dan progress report. & Modalitas: Praktik Lapangan dan Bimbingan. Metode: praktik industri, monitoring berkala, dan bimbingan pembimbing akademik/lapangan. & 1 x 1 x 50' monitoring; kerja mandiri. & Mahasiswa menyerahkan logbook bulan pertama dan melaporkan progress kerja kepada pembimbing akademik. & Bentuk penilaian: Logbook bulan 1 dan progress report. Mengacu rubrik RTM. & Ketepatan waktu; kualitas dokumen; profesionalisme. & 3 \\
5 & SCPMK0804-05703 & Konsultasi pembimbing akademik: identifikasi proyek industri. & Modalitas: Praktik Lapangan dan Bimbingan. Metode: praktik industri, monitoring berkala, dan bimbingan pembimbing akademik/lapangan. & 1 x 1 x 50' konsultasi; persiapan mandiri. & Mahasiswa berkonsultasi dengan pembimbing akademik untuk mengidentifikasi proyek/tugas industri yang relevan dengan kompetensi RPL. & Bentuk penilaian: Notulensi konsultasi dan rencana proyek. Mengacu rubrik RTM. & Ketepatan waktu; kualitas dokumen; profesionalisme. & 3 \\
6 & SCPMK0607-05704 & Monitoring bulan 2: evaluasi penyesuaian dan kompetensi. & Modalitas: Praktik Lapangan dan Bimbingan. Metode: praktik industri, monitoring berkala, dan bimbingan pembimbing akademik/lapangan. & 1 x 1 x 50' monitoring; kerja mandiri. & Mahasiswa dievaluasi kemampuan penyesuaian di lingkungan industri dan progress penerapan kompetensi rekayasa perangkat lunak. & Bentuk penilaian: Logbook bulan 2 dan evaluasi kompetensi. Mengacu rubrik RTM. & Ketepatan waktu; kualitas dokumen; profesionalisme. & 3 \\
7 & SCPMK0804-05703, SCPMK0607-05704 & Mid-internship review bersama pembimbing lapangan. & Modalitas: Praktik Lapangan dan Bimbingan. Metode: praktik industri, monitoring berkala, dan bimbingan pembimbing akademik/lapangan. & 1 x 2 x 50' review; kerja mandiri. & Mahasiswa menjalani review tengah magang bersama pembimbing lapangan untuk mengevaluasi kinerja dan kompetensi yang telah dicapai. & Bentuk penilaian: Mid-internship review report. Mengacu rubrik RTM. & Ketepatan waktu; kualitas dokumen; profesionalisme. & 10 \\
8 & SCPMK0804-05703 & Penilaian pembimbing lapangan mid-term. & Modalitas: Praktik Lapangan dan Bimbingan. Metode: praktik industri, monitoring berkala, dan bimbingan pembimbing akademik/lapangan. & Disesuaikan jadwal industri. & Pembimbing lapangan memberikan penilaian resmi terhadap kinerja, profesionalisme, dan penyelesaian tugas mahasiswa hingga pertengahan magang. & Bentuk penilaian: Formulir penilaian pembimbing lapangan mid-term. Mengacu rubrik RTM. & Ketepatan waktu; kualitas dokumen; profesionalisme. & 15 \\
9 & SCPMK0301-05701 & Monitoring bulan 3: progress laporan magang. & Modalitas: Praktik Lapangan dan Bimbingan. Metode: praktik industri, monitoring berkala, dan bimbingan pembimbing akademik/lapangan. & 1 x 1 x 50' monitoring; penulisan mandiri. & Mahasiswa melaporkan progress penulisan laporan magang dan mendapatkan masukan dari pembimbing akademik. & Bentuk penilaian: Draft awal laporan magang. Mengacu rubrik RTM. & Ketepatan waktu; kualitas dokumen; profesionalisme. & 3 \\
10 & SCPMK0303-05702 & Bimbingan penulisan laporan: struktur dan konten. & Modalitas: Praktik Lapangan dan Bimbingan. Metode: praktik industri, monitoring berkala, dan bimbingan pembimbing akademik/lapangan. & 1 x 1 x 50' bimbingan; penulisan mandiri. & Mahasiswa mendapatkan bimbingan teknis penulisan laporan magang meliputi sistematika, konten, dan tata tulis ilmiah. & Bentuk penilaian: Draft laporan dengan catatan bimbingan. Mengacu rubrik RTM. & Ketepatan waktu; kualitas dokumen; profesionalisme. & 3 \\
11 & SCPMK0303-05702 & Draft laporan magang: bab I--III. & Modalitas: Praktik Lapangan dan Bimbingan. Metode: praktik industri, monitoring berkala, dan bimbingan pembimbing akademik/lapangan. & 1 x 1 x 50' bimbingan; penulisan mandiri. & Mahasiswa menyelesaikan draft BAB I (Pendahuluan), BAB II (Profil Perusahaan), dan BAB III (Pelaksanaan Magang). & Bentuk penilaian: Draft BAB I--III laporan magang. Mengacu rubrik RTM. & Ketepatan waktu; kualitas dokumen; profesionalisme. & 3 \\
12 & SCPMK0607-05704 & Monitoring bulan 4: finalisasi proyek industri. & Modalitas: Praktik Lapangan dan Bimbingan. Metode: praktik industri, monitoring berkala, dan bimbingan pembimbing akademik/lapangan. & Disesuaikan jadwal magang. & Mahasiswa menyelesaikan proyek/tugas industri utama dan mendokumentasikan hasilnya sebagai bahan laporan. & Bentuk penilaian: Laporan finalisasi proyek industri. Mengacu rubrik RTM. & Ketepatan waktu; kualitas dokumen; profesionalisme. & 3 \\
13 & SCPMK0607-05704 & Penilaian pembimbing lapangan akhir. & Modalitas: Praktik Lapangan dan Bimbingan. Metode: praktik industri, monitoring berkala, dan bimbingan pembimbing akademik/lapangan. & Disesuaikan jadwal industri. & Pembimbing lapangan memberikan penilaian akhir atas keseluruhan kinerja dan penyelesaian tugas mahasiswa selama magang. & Bentuk penilaian: Formulir penilaian pembimbing lapangan akhir. Mengacu rubrik RTM. & Ketepatan waktu; kualitas dokumen; profesionalisme. & 15 \\
14 & SCPMK0303-05702 & Revisi dan finalisasi laporan magang. & Modalitas: Praktik Lapangan dan Bimbingan. Metode: praktik industri, monitoring berkala, dan bimbingan pembimbing akademik/lapangan. & 1 x 1 x 50' bimbingan; penulisan mandiri. & Mahasiswa merevisi laporan magang berdasarkan masukan pembimbing akademik dan melengkapi seluruh bab serta lampiran. & Bentuk penilaian: Laporan magang final. Mengacu rubrik RTM. & Ketepatan waktu; kualitas dokumen; profesionalisme. & 5 \\
15 & SCPMK0303-05702 & Pengumpulan dokumen dan persiapan sidang. & Modalitas: Praktik Lapangan dan Bimbingan. Metode: praktik industri, monitoring berkala, dan bimbingan pembimbing akademik/lapangan. & 1 x 1 x 50' persiapan; mandiri. & Mahasiswa mengumpulkan seluruh dokumen persyaratan sidang dan mempersiapkan bahan presentasi sidang magang. & Bentuk penilaian: Kelengkapan dokumen sidang. Mengacu rubrik RTM. & Ketepatan waktu; kualitas dokumen; profesionalisme. & 3 \\
16 & SCPMK0301-05701 & Logbook final dan evaluasi keseluruhan. & Modalitas: Praktik Lapangan dan Bimbingan. Metode: praktik industri, monitoring berkala, dan bimbingan pembimbing akademik/lapangan. & 1 x 1 x 50' evaluasi; mandiri. & Mahasiswa melengkapi dan mengumpulkan logbook final yang mencakup seluruh kegiatan magang, divalidasi pembimbing lapangan. & Bentuk penilaian: Logbook final lengkap. Mengacu rubrik RTM. & Ketepatan waktu; kualitas dokumen; profesionalisme. & 5 \\
17 & SCPMK0303-05702, SCPMK0607-05704 & Sidang magang: presentasi dan defense laporan. & Modalitas: Praktik Lapangan dan Bimbingan. Metode: praktik industri, monitoring berkala, dan bimbingan pembimbing akademik/lapangan. & 1 x 2 x 50' sidang. & Mahasiswa mempresentasikan dan mempertahankan laporan magang di hadapan tim penguji. & Bentuk penilaian: Sidang magang (presentasi dan tanya jawab). Mengacu rubrik RTM. & Ketepatan waktu; kualitas dokumen; profesionalisme. & 20 \\
\end{longtblr}
